% section: おまけ：定数数列

\section{定数数列}\label{節：定数数列}
すべての項が同じ値になる\bookterm{sequence}{数列}を,定数\bookterm{sequence}{数列}という.
元の項が変化していても,作った量は一定になることがある.先ほどの階比型では,
\[
  b_n=\frac{a_n}{n(n+1)},\qquad b_{n+1}=b_n=b_1=1
\]
なので,$a_n=n(n+1)$と元へ戻せた.
一定になる量が見つかると,初項で決まった値がその後も保たれる.
このことが,元の項を求める関係式になる.
また,\sectionref{節：特性方程式型}では,$a_{n+1}=pa_n+q$自身の定数解$a_n=\alpha$を探す.
一定になる値を知ると,そこからの差を追えるようになる.
